\chapter{The Second Page and Direct Second Differentials}
Write $E_2^{s,t}(o)$ for the mod-two Adams second page of a spectrum
$o$, with stem $t-s$. An omitted spectrum means $S^0$.
The permanence argument uses three kinds of information from this page:
small groups near stem $126$, multiplication by the Hopf
classes, and the maps associated to the Hopf cofiber. The most
basic groups are
\[
 E_2^{1,64}(S^0)=\mathbb F_2\{h_6\},\qquad
 E_2^{2,128}(S^0)=\mathbb F_2\{h_6^2\}.
\]
The named classes $W,T,U$ and their companions below describe the
possible obstructions to the survival of $h_6^2$. Product relations and
nondivisibility will eliminate some of these obstructions; four second
differentials supply starting points for the later-page calculations in
Chapter~5.

As in \cite{LWX126}, spectra are $2$-completed, connective and of
finite type. The calculations use the conventions of \cite{LinProgram}.
We first describe a polynomial quotient and its complete bases through
$t\leq261$ for the sphere, then the bases through $t\leq200$ for the Hopf
cofiber $C\nu$. The comparison assumptions identify these models with
Adams groups and specify their products and cofiber maps. For the
top-cell class, they include the four desuspension conditions and the
five basis identifications stated below.

The three additional multiplication equations and four direct $d_2$
equations are separate computational inputs. Their consequences are
proved after each input is stated. The basis notation introduced here
is also used in the
\href{paper/finite-calculations.html}{finite calculation table}
for Chapter~5. The final section gives the numerical presentation and
references.

\section{Algebraic model and sphere comparison}

Monomials in the quotient algebra provide coordinates for the Adams
groups. The complete-basis assumption specifies those coordinates;
the following comparison specifies how multiplication is transported
to the sphere.

\begin{definition}[A finite presentation and complete bases]
\label{sem:Computation:fixed-archive}
\level{coarse}
Let $J$ index the finite ordered list of generators in the sphere
calculation of \cite{LinProgram}, and write
$(s_j,t_j)\in\mathbb N^2$ for their assigned bidegrees. Form the
polynomial ring $R=\mathbb F_2[X_j\mid j\in J]$, with
$\deg X_j=(s_j,t_j)$.
Generator indices start at $0$ and refer to this one ordered list.
Using its complete prescribed family of relations $\mathcal R$, form
\[
 A_{\mathrm{src}}=R/(\mathcal R),\qquad
 \pi:R\longrightarrow A_{\mathrm{src}}.
\]
The generator names used below refer to the residue classes of the
corresponding variables in this presentation, with the same names,
degrees and values as in the calculation.

For $s,t\in\mathbb N$, let $R_{s,t}$ be the span of monomials
$\prod_jX_j^{a_j}$ with $\sum_j a_j(s_j,t_j)=(s,t)$ and put
\[
 P^{s,t}=\pi(R_{s,t})\subseteq A_{\mathrm{src}}.
\]
All relations enter the quotient algebra. The bound below specifies
the range of its complete bases and comparison with Adams groups. Multiplication and the unit are
induced by those of $R$.

For each bidegree, let $N_{s,t}$ be the number of monomials in the
ordered basis list supplied by the calculation. Assume that their
values in $P^{s,t}$ form a complete basis for every $s,t\in\mathbb N$
with $t\le261$: there is a $\mathbb Z$-linear isomorphism
\[
 P^{s,t}\cong\bigoplus_{i<N_{s,t}}\mathbb F_2
\]
whose inverse sends the $i$th coordinate vector to the value of the
$i$th listed monomial, with indices $0\leq i<N_{s,t}$.
Every monomial is evaluated with the same generator values.
\lean{KIP126.Computation.Standard.generators,KIP126.Computation.Standard.atomIndex,KIP126.Computation.Standard.atom_record,KIP126.Computation.Standard.basisCount,KIP126.Computation.Standard.basisMonomial,KIP126.Computation.Standard.archive,KIP126.Computation.SphereArchive.atom_degree,KIP126.Computation.Standard.sourceData,KIP126.Computation.Standard.sourceGenerator,KIP126.Computation.Standard.sourceCoordinates,KIP126.Computation.Standard.sourceBasis_values,KIP126.Computation.Standard.sourceCoordinateConstruction}
\end{definition}

\begin{proposition}[Comparison with the sphere through internal degree 261]
\label{sem:Computation:bounded-comparison}
\level{coarse}
\uses{sem:Computation:fixed-archive,sem:AdamsSynthetic:families,sem:AdamsSynthetic:standard-classes,sem:AdamsSynthetic:sphere-second-page-algebra}
We assume that the algebraic model of
Definition~\ref{sem:Computation:fixed-archive} is identified with the
sphere second page by $\mathbb Z$-linear isomorphisms
\[
c_{s,t}:P^{s,t}\cong_{\mathbb Z}E_2^{s,t}(S^0)
\qquad(s,t\in\mathbb N,\ t\leq261)
\]
such that, for all $s,t,s',t'\in\mathbb N$ with $t+t'\leq261$ and
$x\in P^{s,t}$, $y\in P^{s',t'}$, the source and sphere products satisfy
\[
c_{s+s',t+t'}(xy)=c_{s,t}(x)c_{s',t'}(y).
\]
\lean{KIP126.Computation.Standard.comparison,KIP126.Computation.Standard.comparison_product,KIP126.Computation.Standard.sphereComparisonConstruction}
\end{proposition}

For $0\le i<N_{s,t}$, write $b_{S^0;s,t,i}$ for the image under
$c_{s,t}$ of the $i$th basis monomial. Thus $N_{s,t}$ is the
dimension of the sphere group in the stated range, and the $b_{S^0;s,t,i}$
are its complete ordered basis. The last index is local to the
bidegree $(s,t)$: it numbers basis vectors, whereas $j$ in $X_j$
numbers generators. We now identify the named classes and use these
coordinates for the product calculations.

\section{Sphere classes and multiplication}

The named generators
$h_i,g,e_0,\Delta h_1g,\Delta h_6g,Md_0,[B_4],b,x_{n,s,j}$ have the
bidegrees given below. The notation $x_{n,s,j}$ denotes a named generator in
bidegree $(s,n+s)$: $n$ is its stem, $s$ its Adams filtration, and the
optional third index distinguishes names in the same bidegree. An omitted
third index specifies one named generator.

The names $\Delta h_1g$, $\Delta h_6g$, $Md_0$ and $[B_4]$ denote whole
generator names; $\Delta$ and $M$ are parts of those names. The letter
$b$ abbreviates the single generator label $(\Delta e_1+C_0+h_0^6h_5^2)$,
including the displayed $\Delta e_1$ and $C_0$ as parts of that label.
Some useful generator degrees, including the low classes used in the
extension examples, are
\[
\begin{array}{c|c|c|c}
\text{label}&(s,t)&\text{label}&(s,t)\\\hline
d_0&(4,18)&p&(4,37)\\
g&(4,24)&e_0&(4,21)\\
\Delta h_1g&(9,54)&\Delta h_6g&(9,116)\\
Md_0&(10,69)&[B_4]&(9,69)\\
b&(8,70)&&
\end{array}
\]
\begin{theorem}[The Hopf classes in the algebraic model]
\label{sem:Computation:standard-atoms}
\label{sem:Computation:h6-square-calibration}
\label{sem:Computation:h6-line}
\level{coarse}
\uses{sem:Computation:bounded-comparison,sem:AdamsSynthetic:standard-classes}
For $0\le i\le7$, let $a_i$ be the source generator in bidegree
$(1,2^i)$ corresponding to the Hopf class of
Definition~\ref{sem:AdamsSynthetic:hi-data}. Then
$c(a_i)=h_i$ and $c(1)=1$. Henceforth a generator name also denotes its
image on the sphere second page; in particular this convention applies
to $d_0,p,b$ and the generators listed above.

Products add bidegrees, and powers such as $h_0^2,h_0^6,h_5^2$ are
iterated products. In particular, the source square $a_6^2$ maps to the
class $h_6^2\in E_2^{2,128}$. The first-line component at $(1,64)$
has precisely the two elements
\[
 x\in E_2^{1,64}(S^0)\quad\Longrightarrow\quad x=0\text{ or }x=h_6.
\]
\lean{KIP126.Computation.SpherePresentation,KIP126.Computation.Standard.comparison_unit,KIP126.Computation.Standard.atom,KIP126.Computation.Standard.comparison_h0,KIP126.Computation.Standard.comparison_h1,KIP126.Computation.Standard.comparison_h2,KIP126.Computation.Standard.comparison_h4,KIP126.Computation.Standard.comparison_h5,KIP126.Computation.Standard.comparison_h6,KIP126.Computation.Standard.comparison_h3,KIP126.Computation.Standard.comparison_h7,KIP126.Computation.Standard.sourceHi,KIP126.Computation.Standard.comparison_hi,KIP126.Computation.Standard.atom_h0,KIP126.Computation.Standard.atom_h1,KIP126.Computation.Standard.atom_h2,KIP126.Computation.Standard.atom_h4,KIP126.Computation.Standard.atom_h5,KIP126.Computation.Standard.atom_h6,KIP126.Computation.Standard.h6_square,KIP126.Computation.SpherePresentation.h6_square,KIP126.Computation.Standard.h6_exhaustive}
\end{theorem}
\begin{proof}
The basis of $P^{1,2^i}$ consists of $a_i$. Its image and $h_i$ are
therefore the unique nonzero element of the one-dimensional group
$E_2^{1,2^i}$. Multiplicativity identifies their products, including
$a_6^2$ with $h_6^2$. If $u$ is the preimage of the sphere unit, then
$c(1)c(u)=c(u)=1$, so $c(1)=1$. Finally $N_{1,64}=1$, which gives the
assertion about $h_6$.
\end{proof}

\section{Classes near the Kervaire stem}


\begin{definition}[Classes used in the permanence argument]
\label{sem:Computation:w-t-u}
\level{coarse}
\uses{sem:Computation:standard-atoms}
The following names abbreviate the combinations of sphere classes that
occur in the differentials and extension arguments. First set
\[W=x_{126,8,4}+x_{126,8},\quad T=h_1h_4x_{109,12},\quad U=h_0^2x_{124,8}.\]
Their Adams bidegrees are $(8,134)$, $(14,139)$ and $(10,134)$, respectively.

The companion classes are
\[V=x_{123,9}+h_0x_{123,8},\quad X=h_1x_{121,7},\quad Y=h_0^2x_{125,9,2}.\]
Their bidegrees are $(9,132)$, $(8,130)$ and $(11,136)$. The combination $h_1V+U$ lies in bidegree $(10,134)$.

For the product and boundary calculations, put
$P=h_6Md_0\in E_2^{11,133}$, $Q=h_5x_{91,11}\in E_2^{12,134}$ and $C=e_0\Delta h_6g\in E_2^{13,137}$. Also set $H=g^4\Delta h_1g\in E_2^{25,150}$.

Finally put
\[B=h_5x_{94,8},\qquad B'=h_5x_{94,8}+h_6b\in E_2^{9,134}.\]
The third-differential alternative involving them is stated in
Theorem~\ref{sem:Computation:d3-candidates}.
\lean{KIP126.Computation.Standard.W,KIP126.Computation.Standard.T,KIP126.Computation.Standard.U,KIP126.Computation.Standard.V,KIP126.Computation.Standard.X,KIP126.Computation.Standard.Y,KIP126.Computation.Standard.d2Target,KIP126.Computation.Standard.P,KIP126.Computation.Standard.Q,KIP126.Computation.Standard.correction,KIP126.Computation.Standard.highClass,KIP126.Computation.Standard.d3Candidate,KIP126.Computation.Standard.d3OtherCandidate}
\end{definition}


\section{The group containing $h_6^2$}
\begin{theorem}[The group containing $h_6^2$]
\label{sem:Computation:square-line}
\level{coarse}
\uses{sem:Computation:bounded-comparison,sem:Computation:h6-square-calibration}
The class $h_6^2$ is nonzero and every $x\in E_2^{2,128}$ satisfies $x=0$ or $x=h_6^2$.
\lean{KIP126.Computation.Standard.square_nonzero,KIP126.Computation.Standard.square_exhaustive}
\end{theorem}
\begin{proof}
\uses{sem:E2Computation:two-line-exhaustive}
The complete basis in bidegree $(2,128)$ consists of $h_6^2$, so this
class is nonzero. The description of the Adams two-line in
Theorem~\ref{sem:E2Computation:two-line-exhaustive} gives another way to
identify the unique nonzero element: it is a product $h_ih_j$ with
$i\leq j$ and $2^i+2^j=128$, forcing $i=j=6$.
\end{proof}

\section{The Hopf cofiber and its top-cell class}

\begin{definition}[A class over the top cell of the Hopf cofiber]
\label{sem:Computation:nu-cofiber-object}
\label{sem:Computation:top-cell-comparison}
\label{sem:Computation:nu-cofiber}
\level{coarse}
\uses{sem:AdamsSynthetic:families,sem:Computation:w-t-u,sem:AdamsSynthetic:selected-convergence,sem:Computation:standard-atoms,sem:Computation:products-vanishing,sem:Literature:normalized-e2-long-exact}
Let $C\nu$ be the cofiber of the Hopf map $\nu:S^3\to S^0$. Use its
Adams spectral sequence, and write $i_*$ for the map induced by the
bottom-cell inclusion $i:S^0\to C\nu$.

Choose $a\in E_2^{8,134}(C\nu)$ lying over
$X=h_1x_{121,7}$ on the top cell, with the following four precise
desuspension conditions; existence follows from exactness as explained below. Let
$\delta:C\nu\to\Sigma S^3$ be the connecting map and
$\sigma_{j,1}:\Sigma S^j\cong S^{j+1}$ the shift-coherence isomorphisms.
There are classes
\[
 a_3\in E_2^{8,133}(S^3),\qquad
 a_2\in E_2^{8,132}(S^2),\qquad
 a_1\in E_2^{8,131}(S^1)
\]
such that $\delta_*a$ desuspends to $a_3$,
$(\sigma_{2,1}^{-1})_*a_3$ desuspends to $a_2$,
$(\sigma_{1,1}^{-1})_*a_2$ desuspends to $a_1$, and $a_1$ desuspends to
$X=h_1x_{121,7}$.

For existence, the sphere basis in bidegree $(8,130)$ consists of $X$.
The relation $Xh_2=0$ from
Theorem~\ref{sem:Computation:products-vanishing}, together with the
first-quotient comparison for the normalized $\nu$ map, makes its
$E_2$ operation vanish on the threefold suspension of $X$.
Exactness at $E_2^{8,133}(S^3)$ in
Theorem~\ref{sem:Literature:normalized-e2-long-exact} then supplies $a$.
The intermediate classes are obtained by suspension and the displayed
shift-coherence isomorphisms. The construction uses exactness and the sphere comparison.

Define
\[
 A_{C\nu}=a+i_*x_{126,8,2}+i_*x_{126,8}.
\]
In the module labels for this two-cell cofiber, $x[0]$ denotes the
bottom-cell class $i_*x$, and $[4]$ marks a chosen lift from the top
cell. These labels specify the chosen bottom-cell classes and top-cell lifts.
In the module presentation of $E_2(C\nu)$, the label for $a$ is
$x_{121,7}h_1[4]$, where $h_1[4]$ has bidegree $(1,6)$.
Proposition~\ref{sem:Computation:cofiber-realization-input}
also requires this particular lift $a$ to be the first basis element
in bidegree $(8,134)$, thereby fixing its basis coordinate as well as
its top-cell image.
\lean{KIP126.Computation.Standard.nuCofiberSequence,KIP126.Computation.Standard.bottom,KIP126.Computation.Standard.cofiberTopCellConstruction,KIP126.Computation.Standard.topCell,KIP126.Computation.Standard.top_comparison,KIP126.Computation.Standard.cofiberSource}
\end{definition}


\begin{theorem}[Agreement with the literature labels]
\label{sem:Computation:literature-labels}
\level{coarse}
\uses{sem:Computation:standard-atoms,sem:Literature:kappa-detection}
The images under $c$ of the generators $g$ and $\Delta h_1g$ are the
sphere classes denoted by the same names in \cite{LWX126}.
\lean{KIP126.Computation.Standard.g_agrees,KIP126.Computation.Standard.delta_agrees}
\end{theorem}




\section{Bases and multiplication formulas}


\begin{proposition}[Compatible bases for the Hopf cofiber]
\label{sem:Computation:cofiber-realization-input}
\level{coarse}
\uses{sem:Computation:nu-cofiber,sem:Computation:bounded-comparison}
For $s,t\in\mathbb N$ with $t\leq200$, let $Q^{s,t}$ be the
$\mathbb F_2$-space with basis the ordered monomial labels for the
cofiber, and let $Q^{s,t}=0$ when that list is empty. We assume
isomorphisms
\[
e_{s,t}:E_2^{s,t}(C\nu)\cong_{\mathbb Z}Q^{s,t}.
\]
For $0\leq i<\dim Q^{s,t}$, write $v_{s,t,i}$ for the $i$th
coordinate unit vector and $u_{s,t,i}=e_{s,t}^{-1}(v_{s,t,i})$.
The sphere basis $b_{S^0;s,t,i}$ is the one defined above.
Require the same family $e$ to be compatible with the chosen top-cell
lift and bottom-cell map through all five identities
\[
\begin{aligned}
e_{8,134}^{-1}(v_{8,134,0})&=a,\\
e_{8,134}^{-1}(v_{8,134,3})&=i_*b_{S^0;8,134,2},&
e_{8,134}^{-1}(v_{8,134,4})&=i_*b_{S^0;8,134,3},\\
e_{11,136}^{-1}(v_{11,136,3})&=i_*b_{S^0;11,136,3},&
e_{14,139}^{-1}(v_{14,139,2})&=i_*b_{S^0;14,139,1}.
\end{aligned}
\]
\lean{KIP126.Computation.NuCofiberArchive.realizationConstruction}
\end{proposition}

\begin{theorem}[Complete bases for the required groups]
\label{sem:Computation:table-bases}
\level{coarse}
\uses{sem:Computation:bounded-comparison,sem:Computation:nu-cofiber,sem:Computation:fixed-archive,sem:Computation:cofiber-realization-input}
A cofiber basis label is a monomial in sphere generators acting on one
cofiber module generator, using the sphere-module structure on
$E_2(C\nu)$. Write $b_{C\nu;s,t,i}=u_{s,t,i}$ in the notation of
Proposition~\ref{sem:Computation:cofiber-realization-input}.
These bases are complete for all $t\leq200$. The calculations below use
$71$ of those bidegrees, eight of which have empty basis.

For a spectrum $o=S^0$ or $C\nu$ and a degree $(s,t)$ recorded in the
\href{paper/finite-calculations.html}{finite calculation table}, let
$b_{o;s,t,i}$ be the specified basis vectors. The table writes a class
as a list $I$ of strictly increasing local basis indices, denoting
\[
\sum_{i\in I}b_{o;s,t,i}.
\]
For example, $[0,2]$ means $b_{o;s,t,0}+b_{o;s,t,2}$, and the empty
list means the zero class. The spectrum and bidegree are specified in
the same row. All indices below lie in the corresponding basis range.
For all $s,t\in\mathbb N$ with $t\leq261$, composition with the inverse
sphere comparison gives the $\mathbb F_2$-linear coordinates
\[
E_2^{s,t}(S^0)\cong\bigoplus_{i<N_{s,t}}\mathbb F_2.
\]
The inverse image of the $i$th unit vector is the comparison image
of the $i$th source monomial; applying $c_{s,t}^{-1}$ recovers that
monomial.
The cofiber bases agree with the top-cell lift and bottom-cell images
as follows:
\[
\begin{aligned}
b_{C\nu;8,134,0}&=a,\\
b_{C\nu;8,134,3}&=i_*b_{S^0;8,134,2},&
b_{C\nu;8,134,4}&=i_*b_{S^0;8,134,3},\\
b_{C\nu;11,136,3}&=i_*b_{S^0;11,136,3},&
b_{C\nu;14,139,2}&=i_*b_{S^0;14,139,1}.
\end{aligned}
\]

Thus, in every listed degree, the coordinate isomorphism sends the
displayed basis to the standard basis of an $\mathbb F_2$-vector
space. For the sphere its dimension is $N_{s,t}$ throughout
$t\leq261$; for $C\nu$ the assertion holds throughout $t\leq200$.
\lean{KIP126.Computation.Table.basis_count,KIP126.Computation.Table.cofiber_basis_correct,KIP126.Computation.Table.basis_correct,KIP126.Computation.Table.archive_dimension,KIP126.Computation.Standard.coordinates,KIP126.Computation.Standard.coordinates_basis}
\end{theorem}
\begin{proof}
\uses{sem:Computation:fixed-archive,sem:Computation:bounded-comparison,sem:Computation:nu-cofiber,sem:Computation:cofiber-realization-input}
For the sphere, compose the complete source coordinates with the inverse
of the bounded comparison. The coordinate vectors then correspond to
the listed monomials. For the cofiber, use the equivalences
$e_{s,t}$ for $t\leq200$ and their five compatibility identities from
Proposition~\ref{sem:Computation:cofiber-realization-input}. Their inverse
images of coordinate vectors give the claimed bases, including the empty
bases. Restricting to the selected bidegrees proves the table statements.
The top-cell and connecting-map conditions are retained in this choice
of cofiber bases.
\end{proof}


\begin{theorem}[Named basis vectors and products]
\label{sem:Computation:named-basis-identities}
\level{coarse}
\uses{sem:Computation:table-bases,sem:Computation:standard-atoms,sem:Computation:bounded-comparison,sem:Computation:w-t-u}
For a generator $j\in J$ of bidegree $(s,t)$ with $t\leq261$,
write $G_{j;s,t}$ for its comparison image in $E_2^{s,t}$.
Here $j$ is the generator index and $(s,t)$ records its bidegree.
The monomial expressions defining the sphere basis are evaluated
using these generator images and the sphere unit. Each expression in
position $i$ equals $b_{S^0;s,t,i}$ in the stated bidegree.

The sphere basis vectors satisfy
\[b_{S^0;11,133,1}=P,\qquad b_{S^0;12,134,0}=Q,\qquad b_{S^0;14,139,1}=T,\]
\[b_{S^0;10,134,4}=U,\qquad b_{S^0;13,137,1}=C,\qquad b_{S^0;13,137,2}=h_4x_{109,12}.\]
If $x\in E_2^{s,t}$, $y\in E_2^{s\prime,t\prime}$ and $t+t\prime\leq261$, then $xy=yx$. Moreover, $h_0h_1=h_1h_2=0$.
\lean{KIP126.Computation.Table.basis_P,KIP126.Computation.Table.basis_Q,KIP126.Computation.Table.basis_T,KIP126.Computation.Table.basis_U,KIP126.Computation.Table.basis_correction,KIP126.Computation.Table.basis_h4_x109,KIP126.Computation.Table.product_comm,KIP126.Computation.Table.h0_h1_zero,KIP126.Computation.Table.h1_h2_zero}
\end{theorem}
\begin{proof}
\uses{sem:Computation:table-bases,sem:Computation:bounded-comparison}
In the quotient algebra, the listed monomials representing these basis
vectors equal the displayed products. Applying $c$ gives the six
identities. Likewise commutativity of the polynomial algebra and the
relations $h_0h_1=h_1h_2=0$ pass to the sphere under the multiplicative
comparison.
\end{proof}


\section{The Moore spectrum and its bottom-cell class}

\begin{definition}[The completed Moore spectrum and its bottom class]
\label{sem:AdamsComputation:moore-object}
\label{sem:AdamsComputation:moore-scope}
\level{coarse}
\uses{sem:AdamsSynthetic:families,sem:Computation:table-bases,sem:AdamsSynthetic:paper-spectrum-scope,sem:AdamsSynthetic:homology-exact,sem:AdamsSynthetic:paper-spectrum-instances}
Let $S$ be the completed sphere and choose its Moore cofiber triangle
\[
 C_2=\operatorname{cofib}(2:S\to S),\qquad
 S\xrightarrow{2}S\xrightarrow{i}C_2\xrightarrow{q}\Sigma S.
\]
Its Adams second-page maps are
\[
 i_*:E_2^{s,t}(S)\to E_2^{s,t}(C_2),\qquad
 q_*:E_2^{s,t}(C_2)\to E_2^{s,t}(\Sigma S),
\]
and the selected bottom-cell class is
$\bar x=i_*(b_{S;8,132,0})\in E_2^{8,132}(C_2)$.
These maps and the class $\bar x$ are defined on the actual Adams
spectral sequences.

The spectrum $C_2$ is connective and has finite mod-two homology in
every degree. It therefore belongs to the paper range and to its
suspension-stable bounded-below finite-type enlargement.
\lean{KIP126.AdamsComputation.Moore.spectrum,KIP126.AdamsComputation.Moore.bottom,KIP126.AdamsComputation.Moore.top,KIP126.AdamsComputation.Moore.cofiberSequence,KIP126.AdamsComputation.Moore.bottomPageMap,KIP126.AdamsComputation.Moore.topPageMap,KIP126.AdamsComputation.Moore.bottomSource,KIP126.AdamsComputation.Moore.isPaperSpectrum}
\end{definition}
\begin{proof}
In a negative homotopy degree the two adjacent sphere groups in the
cofiber exact sequence vanish, so the Moore group vanishes. In each
mod-two homology degree, the Moore group is an extension of an image
of one finite sphere group by a subgroup of the next, and is finite.
Completion is supplied by the ambient category.
\end{proof}

The construction in Lemma~\ref{sem:Main:sixth-quotient-representative}
uses these cofiber maps, the normalized exact-sequence comparisons and
the sphere calculations. Its computation input on $C_2$ is the equation
$d_4(\bar x)=0$ in Theorem~\ref{sem:SSComputation:ordinary-differential-inputs}.
Theorem~\ref{sem:SSComputation:selected-class-survival} gives
$\bar x\in Z_4^{8,132}(C_2)$. The nonvanishing of $\bar x$ is proved
from the sphere calculation and the normalized exact sequence in the
proof of that representative construction.

\section{Annihilation relations}


\begin{theorem}[Multiplication inputs and obstruction relations]
\label{sem:AdamsComputation:primitive-relation-inputs}
\label{sem:AdamsComputation:primitive-obstruction-relations}
\level{coarse}
\uses{sem:Computation:bounded-comparison,sem:Computation:named-basis-identities,sem:Computation:fixed-archive}
Let $G_{i;s,t}$ denote the generator image of
Theorem~\ref{sem:Computation:named-basis-identities}. Adopt the following
three equations as computational inputs on the sphere second page:
\[
 h_0G_{24;5,36}=0,\qquad
 (h_0G_{3;1,8})G_{335;12,130}=h_0^2G_{425;12,137},\qquad
 h_2Md_0=h_0G_{89;10,72}.
\]
Together with the algebraic model and its comparison, these three
input equations imply
\[
 h_0^5x_{125,9,2}=0,\qquad h_0^6x_{125,8}=0,\qquad
 h_2G_{407;13,135}=0,\qquad h_0^3G_{89;10,72}=0.
\]
All products here use the sphere $E_2$ multiplication.
\lean{KIP126.Computation.Table.primitive_relation_certificates,KIP126.Computation.Table.h0_obstruction_rel0,KIP126.Computation.Table.h0_obstruction_rel1,KIP126.Computation.Table.h0_obstruction_rel2,KIP126.Computation.Table.h0_obstruction_rel3,KIP126.Computation.Table.h2_obstruction_rel0,KIP126.Computation.Table.h2_obstruction_rel1,KIP126.Computation.Table.h2_obstruction_rel2}
\end{theorem}
\begin{proof}
The relations in $A_{\mathrm{src}}$ give the two vanishing powers of
$h_0$ and the annihilation of $G_{407;13,135}$ by $h_2$; apply the
multiplicative comparison. For the last relation, multiply the assumed
equation $h_2Md_0=h_0G_{89;10,72}$ by $h_0^2$ and use
\[
 h_0^2h_2=h_1^3,\qquad h_1^2Md_0=0.
\]
The left side is then zero. Thus the first three conclusions follow
from the quotient algebra alone, while the fourth also uses the
$h_2Md_0$ input.
\end{proof}

\begin{theorem}[Products with $h_5$]
\label{sem:Computation:theta-product-certificates}
\level{coarse}
\uses{sem:Computation:named-basis-identities}
The multiplication relations in the algebraic model give
\[h_5G_{84;7,70}=0,\qquad h_5G_{85;7,70}=0,\qquad h_5h_6=0.\]
\lean{KIP126.Computation.Table.theta_h5_generator84_zero,KIP126.Computation.Table.theta_h5_generator85_zero,KIP126.Computation.Table.theta_h5_h6_zero}
\end{theorem}

\section{Products, divisibility and second-page vanishing}


\begin{theorem}[Product relations for the named classes]
\label{sem:Computation:products-vanishing}
\level{coarse}
\uses{sem:Computation:w-t-u,sem:Computation:bounded-comparison,sem:Computation:fixed-archive,sem:Computation:standard-atoms}
The following relations will remove products of correction terms in
the later differential arguments:
\[h_5^2b=0,\qquad Xh_2=0,\qquad h_1C=0.\]
\lean{KIP126.Computation.Standard.h5Square_b_zero,KIP126.Computation.Standard.x_h2_zero,KIP126.Computation.Standard.h1_correction_zero}
\end{theorem}
Boundary consequences for $Ph_2$ and $Qh_2$ are stated in
Theorem~\ref{sem:Computation:product-boundaries}.

\begin{theorem}[Nondivisibility of T and Y]
\label{sem:Computation:nondivisibility}
\level{coarse}
\uses{sem:Computation:w-t-u,sem:Computation:bounded-comparison}
The possible obstruction $T$ is divisible by neither $h_0$ nor $h_2$,
and $Y$ is not divisible by $h_2$. More precisely, there is no
$a\in E_2^{13,138}$ with $h_0a=T$, no $a\in E_2^{13,135}$ with
$h_2a=T$, and no $a\in E_2^{10,132}$ with $h_2a=Y$.
\lean{KIP126.Computation.Standard.t_not_h0_multiple,KIP126.Computation.Standard.t_not_h2_multiple,KIP126.Computation.Standard.y_not_h2_multiple}
\end{theorem}
\begin{proof}
\uses{sem:Computation:fixed-archive,sem:Computation:bounded-comparison,sem:Computation:table-bases,sem:Computation:named-basis-identities}
\textit{A coefficient functional for $T$.}
Use the polynomial ring $R$, relation ideal $I=(r_1,\ldots,r_{231848})$
and quotient map $\pi:R\to A_{\mathrm{src}}$ from
\ref{sem:Computation:fixed-archive}. Set
\[
m_1=X_1X_7X_{275},\qquad m_2=X_4X_{320},\qquad
m_3=X_7X_{287},\qquad m_4=X_{70}X_{98}.
\]
For a monomial $m$, write $[m]f$ for its coefficient in $f\in R$.
Define the $\mathbb F_2$-linear functional
\[
L:R\longrightarrow\mathbb F_2,\qquad L(f)=\sum_{k=1}^{4}[m_k]f.
\]
For every defining relation $r_j$, the finite polynomial calculation gives
\[
L(r_j)=0,\qquad [d]r_j=0
\quad\text{for every proper monomial divisor $d$ of any $m_k$.}
\]
The assertion is over the entire list of defining relations, and the
proper divisors include $1$.

To show that $L$ vanishes on $I$, first consider a monomial multiple
$u r_j$. For $u=1$ the first equality applies. For $u\ne1$, the
coefficient of $m_k$ in $u r_j$ is zero if $u$ does not divide $m_k$;
otherwise it is $[m_k/u]r_j$, which vanishes by the proper-divisor
equality. Thus $L(u r_j)=0$ for every monomial $u$. Linearity extends
this to arbitrary polynomial multiples and their sums, giving $L(I)=0$.
Consequently $L$ descends to a linear functional
$\overline L:A_{\mathrm{src}}\to\mathbb F_2$.

None of the four monomials contains $X_0$ or $X_2$, so for every
$a\in A_{\mathrm{src}}$,
\[
\overline L(\pi(X_0)a)=\overline L(\pi(X_2)a)=0,
\qquad \overline L(\pi(m_1))=1.
\]
The generators for $h_0$ and $h_2$ are $\pi(X_0)$ and $\pi(X_2)$,
and the source product for $T=h_1h_4x_{109,12}$ is $\pi(m_1)$.
Hence this source product is neither an $h_0$-multiple nor an
$h_2$-multiple in the quotient algebra.

If $h_0a=T$ held in $E_2(S^0)$, pull $a$ back through
$c_{13,138}$. Multiplicativity, the identifications of $h_0,h_2$ and
injectivity of $c_{14,139}$ would give
$\pi(X_0)c_{13,138}^{-1}(a)=\pi(m_1)$, contradicting the displayed
values of $\overline L$. The same argument with $c_{13,135}$ and $X_2$
excludes $h_2a=T$. Every comparison used here is within $t\leq261$.

\textit{The separate coordinate argument for $Y$.}
The complete basis of $E_2^{10,132}$ is empty, so that group is zero.
On the other hand, $Y=h_0^2x_{125,9,2}$ is the basis vector of index $3$
in bidegree $(11,136)$ and is nonzero. Consequently no class in
$E_2^{10,132}$ can multiply by $h_2$ to give $Y$.

\end{proof}


\begin{theorem}[The required vanishing components]
\label{sem:Computation:vanishing-regions}
\level{coarse}
\uses{sem:Computation:bounded-comparison}
In stem $125$, there are no classes in filtrations zero through four:
the groups $E_2^{s,s+125}$ are zero for $0\leq s\leq4$.
\lean{KIP126.Computation.Standard.e2_stem125_low}
\end{theorem}

\section{Second differentials}

\begin{theorem}[Four second differentials and a named consequence]
\label{sem:AdamsComputation:direct-d2-records}
\label{sem:AdamsComputation:direct-d2-h0-six-h6}
\level{coarse}
\uses{sem:Computation:table-bases,sem:SpectralSequence:differential,sem:Computation:standard-atoms,sem:Computation:named-basis-identities,sem:Computation:bounded-comparison}
Assume the following four differential equations in the sphere Adams
spectral sequence, with source and target given in the complete bases
above:
\[
\begin{aligned}
 d_2(b_{S^0;7,70,2})&=b_{S^0;9,71,0},\\
 d_2(b_{S^0;10,136,0})&=b_{S^0;12,137,2},\\
 d_2(b_{S^0;11,137,0})&=b_{S^0;13,138,3},\\
 d_2(b_{S^0;11,137,3})&=b_{S^0;13,138,4}.
\end{aligned}
\]
These are equations of degree $(2,1)$ on the Adams page, including
the identification of their source and target representatives. They
are the four direct $d_2$ inputs used here, with the computational
conventions of \cite{LinProgram}. The propagation inputs in Chapter~5
are additional assumptions.

The first equation, expressed in named classes, is the nonzero differential
\[
 d_2(h_0^6h_6)=h_0b.
\]
\lean{KIP126.Computation.Table.Raw.directD2Claims,KIP126.Computation.Table.direct_d2_result,KIP126.Computation.Standard.d2_h0Six_h6}
\end{theorem}
\begin{proof}
The source of the first equation is the basis vector
$b_{S^0;7,70,2}=h_0^6h_6$. The source-algebra relation
$h_0b=h_0^7h_5^2$ identifies its target with
$b_{S^0;9,71,0}$, which is nonzero by basis completeness. This proves
the named consequence from the first of the four adopted equations,
using only the $E_2$ algebra and its comparison.
\end{proof}

\begin{theorem}[The classical differential on $h_6$]
\label{sem:AdamsComputation:d2-h6-calibration}
\level{coarse}
\uses{sem:Literature:one-line-differentials,sem:Computation:standard-atoms}
From \cite{Adams60}, as recalled in \cite{LWX126}, Introduction.\par
The Hopf classes satisfy the nonzero differential
\[
d_2(h_6)=h_0h_5^2.
\]
This is the classical one-line Adams differential, used independently
of the four computational equations above.
\lean{KIP126.Computation.Standard.d2_h6}
\end{theorem}
\begin{proof}
\uses{sem:Literature:one-line-differentials}
Specialize the one-line Adams differential to $j=6$ and use the
identifications of the Hopf classes.
\end{proof}

\section{Numerical presentation and references}

The source presentation of Definition~\ref{sem:Computation:fixed-archive}
has $2914$ generators, indexed by $J=\{0,\ldots,2913\}$, and $231848$
relations. In this ordering, $d_0,p,b$ are the residue classes of
$X_8,X_{26},X_{82}$, respectively. All $231848$ relations enter
$R/(\mathcal R)$; the bounds $t\leq261$ and $t\leq200$ specify the
ranges of the sphere and cofiber comparisons.

The complete generator catalogue, relations and dimensions are taken
from the finite sphere calculation of \cite{LinProgram}.
The cofiber data through $t\leq200$ contain $1{,}171$ module generators,
$97{,}621$ module relations and $14{,}989$ ordered basis labels in
$5{,}588$ nonempty bidegrees. The complete numerical data are available
in \cite{LWXZenodoCurrent}, with conventions described in
\cite{LinProgram,LWXMachine}.

The algebraic and comparison inputs are stated in
Definition~\ref{sem:Computation:fixed-archive} and
Propositions~\ref{sem:Computation:bounded-comparison}
and~\ref{sem:Computation:cofiber-realization-input}.
The three multiplication equations and four direct second
differentials are the separate inputs of
Theorems~\ref{sem:AdamsComputation:primitive-relation-inputs}
and~\ref{sem:AdamsComputation:direct-d2-records}.
